\unitheading{K9\quad 刚体运动}{unit-k09}
\addcontentsline{toc}{chapter}{K9 刚体运动}

刚体是内部任意两点距离始终不变的质点系统，是对形变可以忽略的真实物体的理想化。仅给出刚体质心位置还不足以确定其状态，还必须给出姿态。三维姿态可由三个参数描述，欧拉角是常用选择之一。

任何真实物体受力后都会产生程度不同的形变；只有当形变量相对物体尺寸足够小、对所研究问题没有显著影响时，才把它理想化为刚体。由于内部质量分布相对刚体不变，质心在刚体上的位置也固定不动。描述质点只需给出一个位置，描述刚体则既要说明代表其整体平移的位置，通常取质心，又要说明物体朝向，因此自由度更多。

\EulerStepsDiagram

欧拉角不是三根固定轴上的三个独立转角，而是按顺序进行的三次转动：先绕固定轴 $z$ 定出两平面的交线（节线）$N$，再绕 $N$ 改变固连轴 $Z$ 的倾角，最后绕 $Z$ 决定刚体自身方位。高中阶段不要求用它作复杂计算，但应知道三维刚体姿态需要三个自由度。

图2.20中每一步都从正转轴的一端朝原点看，正转向在纸面上表现为逆时针。第二图中 $N$ 垂直纸面向外，$xy$、$XY$ 两平面各表现为一条截线；$z$、$Z$ 分别垂直于它们，所以平面倾角与法线夹角同为 $\beta$。第三步才绕已经倾斜的 $Z$ 改变 $X,Y$ 的方位。三次转轴依次为 $z,N,Z$，不是始终固定的三根坐标轴；交换转动次序通常会改变最终姿态。

常见刚体运动包括：

\begin{itemize}
\item 平动：刚体姿态不变，所有点在同一时刻速度和加速度相同，质心运动即可代表整体运动；
\item 定轴转动：刚体绕空间中固定轴线转动，可用角位置、角速度和角加速度描述；
\item 平面平行运动：平动与转动同时存在、转轴方向保持平行，轮的滚动是典型例子。
\end{itemize}

平面平行运动可以分解为选定参考点的平动与刚体绕该点的转动，通常选质心最方便；也可以在某一瞬间改选速度为零的瞬心，把当时的速度场看成绕瞬心的纯转动。两种分解描述的是同一瞬时运动，参考点的选择不同不会改变刚体各点的实际速度。

定轴转动中，轴上各点速度为零，离轴垂直距离为 $r$ 的点速率为 $|\omega|r$，方向沿当地圆周切线。平面运动在某一瞬间的瞬心具有同样作用：任一点速度垂直于该点到瞬心的连线，大小与距离成正比。因此，已知两个不平行的点速度方向时，可分别作速度的垂线，其交点就是瞬心。若各点速度平行且相等，垂线不相交，瞬心可视为在无穷远处，对应纯平动。

瞬心只由当时的速度场决定，可能随时间移动，也不必与质心重合。质心由质量分布决定，瞬心由运动状态决定；图2.94的不同质量和端点速度组合正是为了把两者区分开。求瞬心时不应预先把质心标在杆的几何中点，更不能因为某点在图上居中就假定它的速度为零。

\RigidMotionTypesDiagram
图中左侧两个姿态的对应边始终平行，各点沿相同形状、彼此平移的曲线运动，说明平动不一定是直线运动。右侧 $P$ 到固定轴 $AB$ 的垂距为 $R$；从参考半径转过角位置 $\theta$ 后，$P$ 的速度仍垂直 $OP$，速率为 $|\omega|R$。后半圆用虚线表示，半径 $OP$ 与点 $P$ 处切线属于同一圆平面。

平面刚体在任一瞬间都可视为绕一个速度为零的点转动，这个点称为瞬时转动中心，简称瞬心。它可以在刚体内部、边界或外部；纯平动的瞬心可视为在无穷远处。

\KinematicKeepSpace{14cm}
\begin{origexamplebox}{例2.94\quad 两质点轻杆刚体的质心与瞬心}
质点 $A,B$ 由长 $l$ 的刚性轻杆连接，此刻杆水平，约定有符号速度的正方向均为竖直向上（图中箭头另按各问实际方向画出）。分别求下列情形中质心位置、速度和瞬心位置：\par
（1）$m_A=m_B=m$，$v_A=2v_B=v_0$；\par
（2）$m_A=m_B=m$，$v_A=-v_B=v_0$；\par
（3）$m_A=2m_B=m$，$v_A=-v_B=v_0$。
\par 题图见图~\ref{question-figure-2-94}。
\KQuestionDQ
\end{origexamplebox}
\begin{xesolution}
图中三种情形使用同一长度标度；$G$ 表示质心，$I$ 表示瞬心。取 $A$ 为 $x=0$、$B$ 为 $x=l$：
\RigidBarCasesDiagram
由质心的质量加权定义，坐标和速度分别为
\[
x_G=\frac{m_Bl}{m_A+m_B},\qquad
v_G=\frac{m_Av_A+m_Bv_B}{m_A+m_B}.
\]
以 $A$ 为参考点，把运动分解成 $A$ 的平动和绕 $A$ 的转动。转动部分的竖直速度为 $\omega x$，而在 $B$ 点有 $\omega l=v_B-v_A$，所以沿杆各点的竖直速度随 $x$ 线性变化：
\[
v(x)=v_A+\frac{v_B-v_A}{l}x.
\]
瞬心 $I$ 满足 $v(x_I)=0$。

（1）$x_G=l/2$，$v_G=3v_0/4$ 向上；
\[
x_I=-\frac{lv_A}{v_B-v_A}=2l,
\]
故瞬心在 $B$ 外侧、距 $B$ 为 $l$。

（2）$x_G=l/2$，$v_G=0$，且 $x_I=l/2$。刚体此刻绕质心纯转动，角速度大小 $2v_0/l$，方向为顺时针。

（3）由 $m_A=m$、$m_B=m/2$，得
\[
\boxed{x_G=l/3},\qquad \boxed{v_G=v_0/3\text{，向上}}.
\]
端点速度仍等大反向，所以瞬心仍在 $\boxed{x_I=l/2}$，不与质量分布一起移动。三种情况的结果分别说明瞬心由速度场决定，而质心由质量分布决定。
\end{xesolution}

\KinematicKeepSpace{10cm}
\begin{origexamplebox}{例2.95\quad 滑动乒乓球转为纯滚动}
半径为 $R$ 的球初始质心速度为 $v_0$、向右，同时以角速度 $\omega_0$ 逆时针转动。转为纯滚动前，质心加速度大小为 $a$、向左，角加速度大小为 $\beta$、顺时针。接触点瞬时速度为零后，质心加速度和角加速度均变为零，球保持匀速纯滚动。求最终质心速度的大小和方向。
\par 题图见图~\ref{question-figure-2-95}。
\KQuestionDR
\end{origexamplebox}
\begin{xesolution}
题图把平动方向、转动方向和接触点放在同一幅图中。取向右和逆时针为正：
\RollingBallDiagram
滑动阶段
\[
v(t)=v_0-at,\qquad \omega(t)=\omega_0-\beta t.
\]
球底接触点速度为 $v+R\omega$；纯滚动临界条件是
\[
v_0-at_*+R(\omega_0-\beta t_*)=0.
\]
故
\[
t_*=\frac{v_0+R\omega_0}{a+R\beta},
\]
最终有符号速度
\[
\boxed{v_f=v_0-at_*
=\frac{R(\beta v_0-a\omega_0)}{a+R\beta}}.
\]
因此速率为 $|v_f|$；$\beta v_0>a\omega_0$ 时向右，反之向左，等号时质心恰停下。此时必满足 $v_f=-R\omega_f$，可作方向校验。
这里归零的是 $a_C$ 和角加速度，而不是所有物质点的加速度。匀速纯滚动时，取通过质心、平行于角速度的轴线，距此轴垂直距离为 $r$ 的球上点仍有大小 $\omega_f^2r$、指向此轴的向心加速度。这里的轴不是经过接触点的瞬时转动轴；底部瞬时接触点速度虽为零，加速度也不必为零。只有 $v_f=\omega_f=0$ 的静止情形才全部为零。
\end{xesolution}

\KinematicKeepSpace{18cm}
\begin{origexamplebox}{例2.96\quad 飞刀的质心抛体与绕质心转动}
在一次飞刀表演中，表演者先后甩出两把飞刀：第一把钉在竖直木靶上，到达时刀身水平；第二把的刀尖恰好扎在第一把的木柄上。观察飞刀姿态序列可把运动分解为两部分：质心 $C$ 作抛体运动，飞刀整体同时绕质心 $C$ 匀速转动。要准确命中，必须同时控制质心出手速度的大小和方向，以及绕质心的角速度。\par
飞刀质心 $C$ 与靶点初始等高，水平距离 $s=5.0\,\mathrm m$；刀轴与水平夹角 $\theta=60^\circ$，刀长 $L=25\,\mathrm{cm}$，质心到刀尖距离 $l=10\,\mathrm{cm}$，$g=9.8\,\mathrm{m/s^2}$。\par
（1）质心按 $30^\circ$ 仰角抛出，使刀尖击中靶点且刀身最终水平、刀尖朝向木靶。刀从初始刀尖朝右上方 $60^\circ$ 的姿态逆时针转动，取达到上述姿态的最小正转角。求质心初速率 $v_C$ 和匀速转动角速度 $\omega$。\par
（2）第二把飞刀从同一位置、同一初始姿态出手，刀尖击中第一把刀木柄的远端；碰撞瞬间第二把刀仍为水平姿态，并取从初始姿态转到该姿态的最小正转角。若两次运动的质心最大高度相同，求 $v_C'$、$\omega'$ 分别比 $v_C$、$\omega$ 小百分之几。
\par 题图见图~\ref{question-figure-2-96}。
\KQuestionDS
\end{origexamplebox}
\begin{xesolution}
第一把刀最终水平时，质心在刀尖后方 $l$，所以质心的水平射程是
\KnifeFlightDiagram
\[
d_1=s-l.
\]
同高斜抛射程 $d_1=v_C^2\sin60^\circ/g$，故
\[
\boxed{v_C=\sqrt{\frac{g(s-l)}{\sin60^\circ}}
\approx7.45\,\mathrm{m/s}}.
\]
飞行时间
\[
T=\frac{2v_C\sin30^\circ}{g}=\frac{v_C}{g}\approx0.760\,\mathrm s.
\]
刀从 $60^\circ$ 逆时针转到下一次“刀尖朝向木靶”的水平姿态时，方向角为 $360^\circ$，角位移为 $2\pi-\theta=5\pi/3$。途中 $180^\circ$ 虽然刀轴也水平，但刀尖朝后，不满足命中姿态。因此
\[
\boxed{\omega=\frac{2\pi-\theta}{T}\approx6.89\,\mathrm{rad/s}}.
\]
若允许额外整周，角速度还可加 $2\pi n/T$；题设取最小正值。图中的等时姿态由 $x_C=d_1t/T$、$y_C=v_C\sin30^\circ\,t-gt^2/2$、$\varphi=60^\circ+300^\circ t/T$ 同时决定。

第二把刀的刀尖击中第一把刀木柄远端，且碰撞时刀身水平。下图的局部放大部分表明：从木靶 $Q$ 到撞击点 $P$ 要先退回第一把刀的全长 $L$，第二把刀的质心 $C_2$ 还要再退回 $l$：
\KnifeImpactDiagram
所以第二把刀质心的水平射程唯一确定为
\[
d_2=s-L-l=4.65\,\mathrm m.
\]
最大高度相同意味着初始竖直分速度相同；始末质心仍等高，所以飞行时间仍为 $T$。因此
\[
v_C'=\sqrt{\left(\frac{d_2}{T}\right)^2+\left(v_C\sin30^\circ\right)^2},
\]
并可写成比值
\[
\frac{v_C'}{v_C}=\frac{\cos30^\circ}{d_1}
\sqrt{d_2^2+d_1^2\tan^2 30^\circ}\approx0.9620.
\]
故 $v_C'$ 约小 $\boxed{3.8\%}$。题设已经规定相同的最终水平姿态、最小正转角和相同飞行时间，故
\[
\boxed{\omega'=\omega},
\]
即减小 $\boxed{0\%}$。木柄上的撞击点与碰撞姿态均已在题设中固定，因此结果没有多解。
\end{xesolution}
